% Chapter - Introduction to Bayesian probability ..............................................................

\chapter{Introduction to Bayesian probability}
\label{chapter6}

\epigraph{\textit{Probability statements are just summaries\\of repeated observations.}}{— W. V. Quine}

Bayesian probability plays a central role in modern statistical reasoning, providing a principled framework for quantifying and updating uncertainty. In this approach, probability is not restricted to long-run frequencies but is interpreted more generally as a measure of uncertainty about events, parameters, or hypotheses. This flexibility allows Bayesian methods to address a wide range of inferential problems in a coherent and unified manner \cite{degroot_2012_probability, jaynes_2003_logic}.

\medskip

As discussed in the previous chapter, classical hypothesis testing is primarily concerned with controlling Type~I and Type~II errors. While this framework has been highly influential, it introduces important limitations: inference is reduced to binary decisions, and conclusions depend critically on arbitrary significance thresholds. Moreover, $p$-values do not directly quantify the probability of a hypothesis being true, which often leads to misinterpretation in practice \cite{lehmann_1959_testing, cox_2006_principles}.

\medskip

Attempts to address these issues include adjusting $p$-values (e.g.\ for multiple comparisons) or returning to Fisher’s original interpretation of significance as a measure of evidence rather than a strict decision rule \cite{fisher_1925_statistical, fisher_1935_design}. However, such modifications do not fully resolve concerns about the reliability and interpretation of statistical conclusions.

\medskip

A more fundamental response is to reconsider the epistemological foundations of statistical inference. Instead of viewing statistics as a tool for hypothesis rejection, one may interpret it as a formal mechanism for updating beliefs in light of evidence. This perspective aligns naturally with Bayesian reasoning, in which prior information is combined with observed data to produce posterior conclusions.

\medskip

From a philosophical standpoint, this shift reflects broader debates on the nature of probability and scientific reasoning. The frequentist interpretation, associated with the work of von Mises \cite{vonmises_1928_truth}, defines probability through limiting frequencies of repeatable experiments, whereas the Bayesian and subjectivist interpretations, developed by Ramsey and de Finetti, treat probability as an expression of rational belief under uncertainty \cite{ramsey_1931_foundations, definetti_1974_probability}. These perspectives are discussed extensively in the philosophy of statistics literature \cite{bandyopadhyay_2011_philosophy}.

\medskip

In contemporary practice, statistics is increasingly viewed as a discipline concerned with reasoning, interpretation, and communication. As Spiegelhalter \cite{spiegelhalter_2019_art} emphasizes, \textit{Statistics is the science of learning from data, and of measuring, controlling and communicating uncertainty}. Within this perspective, Bayesian methods provide a natural and transparent way to represent uncertainty, incorporate prior knowledge, and update conclusions as new data become available.

% Independent and conditional events ..................................................................................
\section{Independent and conditional events}

In many situations involving uncertainty, we are interested not only in whether events occur, but in how the occurrence of one event influences another. For instance, observing an outcome may change our expectations about what follows. The language of conditional probability provides a precise way to describe this dependence, forming the basis for reasoning under uncertainty.

A central question is therefore whether two events are related or not. If knowing that one event has occurred does not change our assessment of the other, the events are said to be independent; otherwise, they are dependent. This distinction is fundamental, as it determines whether information can be ignored or must be incorporated into our reasoning. In particular, conditional probability allows us to update probabilities when new information becomes available, a principle that will be essential when we later consider inverse problems and Bayes’ rule.

\medskip

Events $A$ and $B$ are said to be independent if the occurrence of one does not affect the probability of the other:
\begin{equation}
    P(A \cap B) = P(A)P(B).
\end{equation}

Equivalently,
\begin{equation}
    P(A \mid B) = P(A).
\end{equation}


\textbf{Examples of independent events:}

\begin{itemize}
\item \textit{Two coin tosses:} Let $A$ be “first toss is heads” and $B$ be “second toss is heads.” Then
\[
    P(A)=\tfrac{1}{2}, \quad P(B)=\tfrac{1}{2}, \quad P(A\cap B)=\tfrac{1}{4},
\]
and indeed $P(A\cap B)=P(A)P(B)=\tfrac{1}{4}$, so the events are independent.

\item \textit{Two dice:} Let $A$ be “first die shows a 6” and $B$ be “second shows even number.” Then
\[
    P(A)=\tfrac{1}{6}, \quad P(B)=\tfrac{3}{6}=\tfrac{1}{2}, \quad P(A\cap B)=\tfrac{1}{12},
\]
which again satisfies $P(A\cap B)=P(A)P(B)$.
\end{itemize}

When this property fails, events are what we call \textit{dependent}, and conditioning becomes essential:
\begin{equation}
    P(A \mid B)=\frac{P(A\cap B)}{P(B)}.
\end{equation}


\textbf{Examples of dependent events:}

\begin{itemize}
\item \textit{Drawing cards without replacement.} Consider a standard deck of 52 cards. In this case, it we let $A$ be “first card is an Ace” and $B$ be “second card is an Ace.” Then
\[
    P(A)=\tfrac{4}{52}, \qquad P(B)=\tfrac{4}{52},
\]
but we have to be careful,
\[
    P(B \mid A)=\tfrac{3}{51} \neq \tfrac{4}{52},
\]
since one Ace has already been removed from the deck. Hence the events are dependent.

\item \textit{Urn example.} Suppose an urn contains 3 red and 2 blue balls. Let $A$ be “first draw is red” and $B$ be “second draw is red,” without replacement. Then
\[
    P(A)=\tfrac{3}{5}, \qquad P(B \mid A)=\tfrac{2}{4}=\tfrac{1}{2},
\]
while the second ball is again conditional to which one we picked first:
\[
    P(B)=\tfrac{3}{5}.
\]
Thus $P(B \mid A) \neq P(B)$, confirming dependence.
\end{itemize}

% The Bayes rule ......................................................................................................
\section{The Bayes rule}

The notion of conditional probability introduced in the previous section naturally leads to a fundamental question: how can we reason in the reverse direction? That is, if we observe an event $B$, how should we update our belief about a possible cause $A$? Problems of this kind are often referred to as \emph{inverse problems}, and they lie at the heart of statistical inference.

\medskip

Bayes’ rule provides a principled solution to this question. It follows directly from the symmetry of joint probability,
\begin{equation}
    P(A \cap B) = P(B \cap A),
\end{equation}

together with the definition of conditional probability. Rearranging terms yields
\begin{equation}
    P(A \mid B) = \frac{P(B \mid A)P(A)}{P(B)}.
\end{equation}

This formula shows how prior beliefs $P(A)$ are updated in light of observed evidence $B$ through the likelihood $P(B \mid A)$. The elements of this expression are named as follows:

\begin{itemize}
    \item{Prior.}
    \item{Likelihood.}
    \item{Total marginal probability.}
    \item{Posterior, updated, conditional probability.}
\end{itemize}

\medskip

Historically, this idea originates in the work of Thomas Bayes, whose posthumous essay \cite{bayes_1763_doctrine} addressed the problem of inferring causes from observed events. The significance of this result was later recognized and greatly extended by Laplace, who developed it into a general framework for scientific reasoning \cite{laplace_1812_theorie}. Philosophically, Bayes’ rule embodies the idea that learning from data consists in updating prior beliefs in a consistent and quantitative manner.

\medskip

\textbf{Example 1: biased vs fair die.}

Suppose we are given two dice: one fair and one biased such that $P(6)=\tfrac{1}{2}$. One die is chosen at random and rolled once, resulting in a $6$. What is the probability that the chosen die is biased?

\medskip

\textit{Hypotheses:}
\[
H_1:\ \text{die is biased}, \qquad H_2:\ \text{die is fair}.
\]

\textit{Priors:}
\[
P(H_1)=\tfrac{1}{2}, \qquad P(H_2)=\tfrac{1}{2}.
\]

\textit{Likelihoods:}
\[
P(6 \mid H_1)=\tfrac{1}{2}, \qquad P(6 \mid H_2)=\tfrac{1}{6}.
\]

\textit{Marginal probability:}
\[
P(6)=P(6 \mid H_1)P(H_1)+P(6 \mid H_2)P(H_2)
= \tfrac{1}{2}\cdot\tfrac{1}{2}+\tfrac{1}{6}\cdot\tfrac{1}{2}
= \tfrac{1}{4}+\tfrac{1}{12}
= \tfrac{1}{3}.
\]

\textit{Posterior:}
\[
P(H_1 \mid 6)=\frac{P(6 \mid H_1)P(H_1)}{P(6)}
=\frac{\tfrac{1}{2}\cdot\tfrac{1}{2}}{\tfrac{1}{3}}
=\frac{\tfrac{1}{4}}{\tfrac{1}{3}}
=\tfrac{3}{4}.
\]

\medskip

\textbf{Example 2: medical testing.}

Suppose a disease has prevalence $P(D)=0.05$. A test has sensitivity $P(+ \mid D)=0.99$ and false positive rate $P(+ \mid D^c)=0.01$. Given a positive test result, what is the probability of having the disease?

\medskip

\textit{Hypotheses:}
\[
H_1:\ \text{disease present}, \qquad H_2:\ \text{no disease}.
\]

\textit{Priors:}
\[
P(H_1)=0.05, \qquad P(H_2)=0.95.
\]

\textit{Likelihoods:}
\[
P(+ \mid H_1)=0.99, \qquad P(+ \mid H_2)=0.01.
\]

\textit{Marginal probability:}
\[
P(+)=0.99\cdot 0.05 + 0.01\cdot 0.95
=0.0495+0.0095
=0.059.
\]

\textit{Posterior:}
\[
P(H_1 \mid +)=\frac{0.99\cdot 0.05}{0.059}
\approx 0.84.
\]

\medskip

\textbf{Example 3: Monty Hall problem (extended).}

Suppose one of three doors hides a prize. A contestant selects a door. The host, who knows where the prize is, opens one of the remaining doors revealing no prize, and offers the option to switch.

\medskip

\textit{Hypotheses:}
\[
H_1:\ \text{initial choice is correct}, \qquad
H_2:\ \text{initial choice is incorrect}.
\]

\textit{Priors:}
\[
P(H_1)=\tfrac{1}{3}, \qquad P(H_2)=\tfrac{2}{3}.
\]

Let $E$ be the event “host opens a door with no prize.”

\medskip

\textit{Likelihoods:}
\[
P(E \mid H_1)=\tfrac{1}{2}, \qquad P(E \mid H_2)=1,
\]
since if the initial choice is correct, the host has two possible doors to open, whereas if it is incorrect, the host has only one choice.

\medskip

\textit{Marginal probability:}
\[
P(E)=\tfrac{1}{2}\cdot\tfrac{1}{3} + 1\cdot\tfrac{2}{3}
=\tfrac{1}{6}+\tfrac{2}{3}
=\tfrac{5}{6}.
\]

\textit{Posterior:}
\[
P(H_1 \mid E)=\frac{\tfrac{1}{2}\cdot\tfrac{1}{3}}{\tfrac{5}{6}}
=\frac{\tfrac{1}{6}}{\tfrac{5}{6}}
=\tfrac{1}{5}.
\]

Thus,
\[
P(\text{win by switching}) = 1 - P(H_1 \mid E) = \tfrac{4}{5},
\]
illustrating how conditioning on the host’s action changes the probability structure and favors switching.

\medskip

The Monty Hall problem became widely known through its popularization in the late twentieth century, notably in discussions following its presentation in a magazine column \cite{vos_savant_1990_monty}. It is often perceived as paradoxical because it contradicts the intuitive belief that, after one door is opened, the remaining doors should each have probability $\tfrac{1}{2}$. The resolution lies in recognizing that the host’s action is not independent of the underlying state: the information revealed is conditional and therefore alters the probability space. This example highlights the importance of carefully accounting for the data-generating mechanism when applying conditional probability and Bayes’ rule.

% Frequentist vs Bayesian .............................................................................................
\section{Frequentist vs Bayesian}

The frequentist interpretation defines probability as long-run relative frequency and treats parameters as fixed but unknown constants, while randomness arises from sampling variation \cite{vonmises_1928_truth,degroot_2012_probability}. In this view, probability statements are understood through hypothetical repetition: probabilities describe the behavior of a procedure if it were repeated indefinitely under identical conditions.

\medskip

In this framework, statements such as “the probability that the parameter lies in an interval is 0.95” are avoided, since the parameter is not random; instead, one constructs procedures that produce intervals containing the true parameter in 95\% of repeated samples. This leads to the notion of confidence intervals, which are properties of methods rather than probability statements about specific parameters.

\medskip

The Bayesian interpretation, by contrast, treats probability as a coherent degree of belief constrained by rational consistency, allowing probabilities to be assigned to hypotheses and parameters themselves \cite{ramsey_1931_foundations,definetti_1974_probability}. This interpretation extends the use of probability beyond repeatable experiments to situations involving uncertainty about fixed but unknown quantities.

\medskip

Under this view, parameters are modeled as random variables with prior distributions reflecting initial uncertainty, and inference consists of updating this distribution through Bayes’ rule when data are observed. The resulting posterior distribution directly answers questions of interest, such as the probability that a parameter lies within a given interval.

\medskip

Philosophically, the difference concerns meaning rather than mathematics: frequentism anchors probability in hypothetical repetition, whereas Bayesianism interprets probability as rational belief updated by evidence. As a consequence, Bayesian inference provides a unified framework in which learning, uncertainty, and decision-making are naturally integrated within the probabilistic model itself.

% Bayesian inference ..................................................................................................
\section{Bayesian inference}

Bayesian probability provides a general interpretation of probability as a degree of belief under uncertainty. Bayesian \emph{statistics}, in contrast, is the application of this interpretation to infer unknown quantities from data. In this setting, unknown parameters are treated as random variables, and inference consists in updating our beliefs about these parameters after observing data.

More precisely, Bayesian statistics combines prior information with observed evidence in a coherent mathematical framework. This represents a conceptual shift from classical approaches: rather than estimating fixed but unknown parameters, we describe uncertainty about them through probability distributions and update these distributions as new data become available.

\medskip

Bayesian inference is based on Bayes’ rule applied to parameters:
\[
\pi(\theta \mid x)=\frac{p(x\mid \theta)\,\pi(\theta)}{p(x)}.
\]

Each term has a clear interpretation:
\begin{itemize}
\item $\pi(\theta)$ is the \textbf{prior}, representing initial beliefs about the parameter $\theta$,
\item $p(x \mid \theta)$ is the \textbf{likelihood}, describing how probable the observed data $x$ are for a given $\theta$,
\item $\pi(\theta \mid x)$ is the \textbf{posterior}, representing updated beliefs after observing $x$,
\item $p(x)$ is a normalizing constant ensuring the posterior integrates to one.
\end{itemize}

This relationship is often summarized as
\[
\text{Posterior} \propto \text{Likelihood} \times \text{Prior},
\]
emphasizing that learning consists of combining prior information with evidence.

\medskip

\textbf{Example: coin flipping.}

\medskip

Suppose we want to estimate the probability $\theta$ of obtaining heads when tossing a coin. Before observing any data, we express our uncertainty through a prior distribution. A common choice is the Beta distribution:
\[
\theta \sim \mathrm{Beta}(\alpha,\beta),
\]
which is defined on the interval $[0,1]$ and is therefore well-suited for modeling probabilities. Its parameters $\alpha$ and $\beta$ can be interpreted as prior counts: roughly speaking, $\alpha-1$ prior “heads” and $\beta-1$ prior “tails.”

Now suppose we observe $n$ coin tosses, with $k$ heads. The likelihood of the data is binomial:
\[
p(x \mid \theta) \propto \theta^k (1-\theta)^{n-k}.
\]

A key feature of the Beta distribution is that it is \emph{conjugate} to the binomial likelihood, meaning the posterior has the same functional form:
\[
\theta \mid x \sim \mathrm{Beta}(\alpha+k,\beta+n-k).
\]

\textit{Interpretation:} the posterior parameters are obtained by adding the observed counts to the prior counts. For example, if $\alpha=\beta=1$ (a uniform prior) and we observe $k=7$ heads in $n=10$ tosses, then
\[
\theta \mid x \sim \mathrm{Beta}(8,4),
\]
reflecting an updated belief that favors larger values of $\theta$.

\medskip

\textbf{Frequentist comparison (coin example).}  

In a classical hypothesis testing framework, we would approach this problem differently. We would formulate hypotheses such as
\[
H_0:\ \theta = 0.5 \quad \text{(fair coin)}, \qquad
H_1:\ \theta \neq 0.5,
\]
collect data (the $n$ coin tosses), and compute a test statistic, for example the number of heads $k$ or the sample proportion $\hat{\theta}=k/n$. We would then calculate a $p$-value, representing how likely it is to observe data at least as extreme as $k$ under the assumption that $H_0$ is true. Based on this $p$-value, we would decide whether to reject $H_0$ or not. In contrast, the Bayesian approach does not produce a binary decision: instead, it yields a full distribution for $\theta$, directly quantifying uncertainty and allowing probabilistic statements about the parameter.

\medskip

\textbf{Example: Gaussian mean.}

\medskip

Suppose we observe data $X_1,\dots,X_n$ that are normally distributed:
\[
X_i \sim \mathcal{N}(\mu,\sigma^2),
\]
where the variance $\sigma^2$ is known, and the goal is to infer the unknown mean $\mu$.

We begin by expressing our prior belief about $\mu$:
\[
\mu \sim \mathcal{N}(\mu_0,\sigma_0^2),
\]
where $\mu_0$ represents our initial guess for the mean, and $\sigma_0^2$ reflects how uncertain we are about this guess (a larger value means more uncertainty).

\medskip

After observing data $x=(x_1,\dots,x_n)$, the likelihood function describes how plausible different values of $\mu$ are given the data:
\[
p(x \mid \mu) \propto \exp\!\left(-\frac{1}{2\sigma^2}\sum_{i=1}^n (x_i-\mu)^2\right).
\]
This expression is maximized when $\mu$ is close to the sample mean
\[
\bar{x} = \frac{1}{n}\sum_{i=1}^n x_i,
\]
so the data suggest values of $\mu$ near $\bar{x}$.

\medskip

The prior also has a Gaussian form:
\[
\pi(\mu) \propto \exp\!\left(-\frac{1}{2\sigma_0^2}(\mu-\mu_0)^2\right),
\]
which favors values of $\mu$ close to $\mu_0$.

\medskip

Bayesian inference combines these two pieces of information:
\[
\pi(\mu \mid x) \propto p(x \mid \mu)\,\pi(\mu).
\]
Multiplying the two exponentials results in another quadratic expression in $\mu$, which corresponds again to a Gaussian distribution. This is why the normal prior is called \emph{conjugate} to the normal likelihood.

\medskip

The posterior distribution is therefore
\[
\mu \mid x \sim \mathcal{N}\!\left(
\frac{\tfrac{n}{\sigma^2}\bar{x} + \tfrac{1}{\sigma_0^2}\mu_0}{\tfrac{n}{\sigma^2}+\tfrac{1}{\sigma_0^2}},
\;\frac{1}{\tfrac{n}{\sigma^2}+\tfrac{1}{\sigma_0^2}}
\right).
\]

\medskip

\textit{Interpretation.} The posterior mean is a weighted average of two quantities:
\begin{itemize}
\item the sample mean $\bar{x}$ (information from the data),
\item the prior mean $\mu_0$ (initial belief).
\end{itemize}

The weights depend on the \emph{precision} (inverse variance). If we have many observations (large $n$), the term $\tfrac{n}{\sigma^2}$ becomes large, and the posterior mean is close to $\bar{x}$. If instead the prior is very precise (small $\sigma_0^2$), more weight is given to $\mu_0$.

\medskip

In simple terms, the posterior balances prior belief and observed data: it moves our estimate of $\mu$ toward the data, but not completely, unless the data are very informative.

\medskip

\textbf{Frequentist comparison (Gaussian mean).}  

Similarly, in the Gaussian setting, a classical approach would begin by specifying hypotheses such as
\[
H_0:\ \mu = \mu_0, \qquad H_1:\ \mu \neq \mu_0.
\]
After collecting the data, we would compute a test statistic (typically based on the sample mean $\bar{x}$) and standardize it to form a $z$-score. From this, we would derive a $p$-value and decide whether to reject $H_0$ at a chosen significance level. Alternatively, we might construct a confidence interval for $\mu$. However, such intervals are defined in terms of repeated sampling and do not assign probabilities to specific parameter values. By contrast, the Bayesian approach produces a posterior distribution for $\mu$, allowing us to state directly how plausible different values of $\mu$ are given the observed data.
